\chapter{Meningococcal Disease Testing}

\section*{What Is This Test?}
Meningococcal disease testing detects \textit{Neisseria meningitidis} (meningococcus), an aerobic gram-negative diplococcus that causes devastating invasive disease including meningitis (infection of the meninges lining the brain and spinal cord) and meningococcemia (fulminant sepsis with hemorrhagic skin lesions) \cite{stephens2007epidemic}. \textit{N. meningitidis} is classified into 12 serogroups based on capsular polysaccharide; serogroups A, B, C, W, X, and Y account for virtually all invasive disease \cite{borrow2017global}. Transmission is via respiratory droplets; asymptomatic nasopharyngeal carriage (present in 5--25\% of adolescents and young adults) may precede invasive disease. Clinical manifestations include fever, headache, neck stiffness, photophobia, and altered mental status (meningitis); petechial/purpuric rash that progresses rapidly to purpura fulminans, disseminated intravascular coagulation (DIC), and multi-organ failure (meningococcemia). The case fatality rate is 10--15\% despite appropriate antibiotic therapy, and 11--19\% of survivors have permanent sequelae (hearing loss, neurologic disability, limb loss) \cite{stephens2007epidemic}. The incubation period is 1--10 days and disease onset can be fulminant, with death occurring within 24 hours of symptom onset.

\section*{Alternative Names}
Meningococcus Testing; Neisseria meningitidis Culture; Meningococcal PCR; CSF Meningitis Panel; Meningococcal Serogroup; Meningitis Blood Culture; Meningococcal Antigen

\section*{Principle}
Gram stain and culture: CSF, blood, or petechial lesion aspirate Gram stain shows gram-negative diplococci (kidney-bean shaped) within neutrophils. Cultures on chocolate agar or 5\% sheep blood agar incubated at 35--37\textdegree{}C in 5--10\% CO$_2$. Colonies are oxidase-positive, catalase-positive, and identified by carbohydrate utilization (glucose and maltose positive; lactose, sucrose negative), MALDI-TOF, or latex agglutination. Blood cultures become positive in 50--60\% of cases. Meningococcal serogroup is determined by slide agglutination with serogroup-specific antisera or by PCR. PCR: Multiplex real-time PCR targeting \textit{ctrA} (capsule transport gene), \textit{porA}, and serogroup-specific genes (\textit{sacB} for serogroup A, \textit{synD} for B, etc.) directly from CSF, blood (EDTA), or plasma. Sensitivity 80--100\%, specificity >99\%. Results available in 2--4 hours \cite{brouwer2010epidemiology}. CSF latex agglutination test for capsular polysaccharide (A, B, C, Y, W135) has moderate sensitivity (50--80\%) and is now a supplementary test given availability of PCR. Meningococcal multiplex panels (BioFire FilmArray Meningitis/Encephalitis Panel) include \textit{N. meningitidis} detection alongside 13 other pathogens.

\section*{History}
Epidemic meningococcal disease was first described by Gaspard Vieusseux in 1805 in Geneva. The organism was isolated from CSF by Anton Weichselbaum in 1887, who named it \textit{Diplococcus intracellularis meningitidis}. The role of asymptomatic nasopharyngeal carriage was elucidated by Glover in the 1910s. Meningococcal polysaccharide vaccines (bivalent A/C, quadrivalent A/C/Y/W-135) were developed in the 1970s. Meningococcal conjugate vaccines (MenACWY) were introduced beginning in 2005 (Menactra) and recommended for all adolescents \cite{cohn2013prevention}. Serogroup B vaccines (Bexsero, Trumenba), targeting outer membrane protein antigens rather than capsular polysaccharide, were licensed in 2014--2015 \cite{cohn2013prevention}. The rapid progression and high mortality of meningococcemia have made the disease a paradigm for the importance of early empiric antibiotic therapy and post-exposure chemoprophylaxis. PCR for non-culture detection became widely available in the 2000s and has substantially improved diagnostic yield, particularly in patients who receive antibiotics before specimen collection. The FilmArray ME Panel (FDA-cleared 2015) includes \textit{N. meningitidis} in a rapid multiplex format.

\section*{Why This Test Is Ordered}
\begin{itemize}
  \item Diagnosis of acute bacterial meningitis in patients presenting with fever, headache, neck stiffness, photophobia, and altered mental status, with or without petechial/purpuric rash
  \item Rapid diagnosis of meningococcemia in patients with acute onset of fever, petechial/purpuric rash, hypotension, and DIC (Waterhouse-Friderichsen syndrome)
  \item Evaluation of febrile patients with petechial rash, even in the absence of meningeal signs, to exclude early meningococcal disease
  \item Investigation of suspected meningococcal disease outbreaks in closed populations (college dormitories, military barracks, prisons) to guide chemoprophylaxis decisions for close contacts
  \item Confirmation of presumptive meningococcal diagnosis when antibiotics were administered prior to specimen collection, rendering culture negative (PCR may remain positive)
  \item Identification of the serogroup to guide vaccination strategies: the serogroup B vaccine (Bexsero, Trumenba) is distinct from the quadrivalent conjugate (MenACWY) vaccine and has different indications for post-exposure prophylaxis
\end{itemize}

\section*{Primary vs Secondary Test}
Diagnosis of meningococcal disease is a medical emergency. CSF Gram stain and culture plus blood cultures are primary tests that must be obtained STAT. PCR on CSF and/or blood is a primary, not secondary, diagnostic modality and should be performed simultaneously with culture. The FilmArray ME Panel is increasingly used as a primary rapid diagnostic. Latex agglutination is a secondary supplementary test.

\section*{How This Test Is Performed}
Blood cultures: Two sets (aerobic and anaerobic) collected before antibiotic administration from separate venipuncture sites; 20 mL per set in adults. CSF: Collected by STAT lumbar puncture in sterile tubes. Tube 1 for chemistry (glucose, protein), Tube 2 for Gram stain, culture, and PCR, Tube 3 for cell count with differential, Tube 4 for additional studies. Specimen should be transported immediately and not refrigerated (meningococci are temperature-sensitive). If lumbar puncture is contraindicated (coagulopathy, intracranial mass effect, hemodynamic instability), blood cultures and PCR should be obtained and empiric antibiotics administered without delay. Lymph node aspiration or biopsy may be cultured if present. Petechial lesion scraping: The petechial lesion is scraped with a scalpel blade, smeared on a glass slide, Gram stained for gram-negative diplococci (sensitivity 50--70\%).

\section*{What Preparation Is Needed}
No special patient preparation. Urgent collection and transport are paramount. Antibiotics (ceftriaxone, cefotaxime, or penicillin G) should be administered as soon as the clinical diagnosis is suspected and blood cultures drawn; do NOT delay antibiotic therapy to obtain CSF if lumbar puncture cannot be performed immediately. Prehospital administration of parenteral antibiotics (by EMS or referring hospital) is recommended when meningococcal disease is suspected.

\section*{Contraindications and Precautions}
Lumbar puncture is contraindicated in the presence of coagulopathy (INR >1.4, platelets <50,000/$\mu$L), suspected intracranial mass lesion with mass effect (papilledema, focal neurologic deficit, recent seizure), or infection at the LP site \cite{tunkel2004practice}. In these circumstances, blood cultures and PCR should be obtained, antibiotics started, and LP deferred until safe. Prior antibiotic therapy significantly reduces culture sensitivity but has less effect on PCR. Antibiotic treatment should not be withheld for diagnostic testing. Contact and droplet precautions should be implemented for hospitalized patients with suspected meningococcal disease until 24 hours of effective antibiotic therapy is completed, as transmission to healthcare workers has occurred.

\section*{Risks}
The primary risk is failure to diagnose meningococcal disease promptly; mortality reaches 80--100\% without treatment. LP carries risks of post-lumbar puncture headache (10--30\%), bleeding, infection, or brain herniation in the presence of increased intracranial pressure. These risks are minimal when the procedure is performed with appropriate precautions and contraindications are respected. Healthcare workers and laboratory personnel handling specimens from patients with suspected meningococcal disease should use standard precautions and avoid sniffing culture plates, as laboratory-acquired meningococcal disease has occurred.

\section*{What Happens After the Test}
CSF Gram stain results: STAT, within 1 hour. Blood culture: Preliminary Gram stain at 12--24 hours; final identification and sensitivity at 48--72 hours. PCR: 2--4 hours (conventional); 1 hour (FilmArray ME Panel). Positive Gram stain or PCR results are immediately reported as a critical value. Public health authorities must be notified of all confirmed and suspected cases of meningococcal disease (nationally notifiable condition). Close contacts (household members, daycare contacts, anyone exposed to the patient's oral secretions in the 7 days before illness onset) require chemoprophylaxis with rifampin (600 mg PO BID x 2 days or 10 mg/kg BID x 2 days for children), ciprofloxacin (500 mg single dose), or ceftriaxone (250 mg IM single dose) ideally within 24 hours of exposure identification. The index case should receive chemoprophylaxis prior to discharge unless treated with cefotaxime or ceftriaxone, which eradicate nasopharyngeal carriage.

\section*{Understanding Results}

\subsection*{Below Normal}
\begin{itemize}
  \item \textbf{Gram stain negative, culture negative, PCR negative:} No evidence of meningococcal infection. However, cultures may be falsely negative if antibiotics were administered prior to specimen collection (up to 50\% reduction in sensitivity for CSF and blood). A negative result by all three methods in a patient who has not received antibiotics makes meningococcal disease very unlikely
  \item \textbf{CSF profile with low WBC but no organisms:} Very early in meningococcal infection (<12 hours), CSF may be acellular or have few WBCs but still contain viable organisms (fulminant meningococcemia). A normal CSF cell count does not exclude meningococcal meningitis; Gram stain and culture remain essential
\end{itemize}

\subsection*{Above Normal}
\begin{itemize}
  \item \textbf{Gram stain positive for gram-negative diplococci in CSF, blood, or petechial scraping:} Presumptive diagnosis of meningococcal infection warranting immediate continuation of empiric antibiotics (ceftriaxone or cefotaxime). The presence of organisms within polymorphonuclear leukocytes (intracellular) confirms true infection rather than contamination
  \item \textbf{Culture positive for \textit{Neisseria meningitidis}:} Definitive diagnosis. Serogroup identification is essential for public health and chemoprophylaxis decisions. Antimicrobial susceptibility testing for penicillin, ceftriaxone, ciprofloxacin, and rifampin should be performed; penicillin-resistant meningococci remain uncommon in the United States but are increasing globally
  \item \textbf{Meningococcal PCR positive (CSF or blood):} Diagnostic of meningococcal infection even when culture is negative. Includes serogroup identification. PCR is particularly valuable when antibiotics were administered before specimen collection; it can detect nonviable organisms
  \item \textbf{FilmArray ME Panel positive for \textit{N. meningitidis}:} Rapid diagnostic result from multiplex PCR; should be acted upon empirically and confirmed by culture/serogroup testing if possible. The panel does not provide serogroup information
\end{itemize}

\section*{Normal Results}
CSF: \textbf{WBC 0--5 cells/$\mu$L} (adults; up to 30 in neonates), \textbf{glucose >40 mg/dL, protein <45 mg/dL. Gram stain: No organisms seen.} Culture: \textbf{No growth at 5 days.} PCR: \textbf{\textit{N. meningitidis} not detected.}

\section*{Follow-up Tests}
\begin{itemize}
  \item \textbf{Meningococcal serogroup determination:} Slide agglutination or PCR-based serogrouping (A, B, C, W, X, Y) from culture isolates or directly from clinical specimens; critical for guiding vaccination and chemoprophylaxis policies \cite{cohn2013prevention}
  \item \textbf{Antimicrobial susceptibility testing (AST):} Broth microdilution or Etest for penicillin, ampicillin, ceftriaxone, cefotaxime, rifampin, ciprofloxacin, meropenem; all clinical isolates from sterile sites should undergo AST
  \item \textbf{CSF analysis parameters (cell count with differential, glucose, protein):} Typical bacterial meningitis profile: WBC 1,000--5,000 (predominantly neutrophils), glucose <40 mg/dL (CSF:serum glucose ratio <0.4), protein >100 mg/dL. Serial monitoring may guide response to therapy
  \item \textbf{Complete blood count with differential:} Leukocytosis with left shift and toxic granulation; thrombocytopenia due to DIC in fulminant meningococcemia
  \item \textbf{Coagulation studies (PT, PTT, fibrinogen, D-dimer):} DIC screen is indicated in patients with purpura fulminans or petechial rash suggesting active DIC
  \item \textbf{Close contact evaluation and nasopharyngeal culture:} Nasopharyngeal cultures are NOT routinely recommended for management of contacts; chemoprophylaxis should be administered based on epidemiologic risk, not culture status
\end{itemize}

\section*{References}
\bibliographystyle{unsrtnat}
\bibliography{references}

\clearpage
